% Topic content; common layout lives in the shared course preamble.
\newcommand{\CourseTemplatePath}{../../../../../templates/slides/}
\input{\CourseTemplatePath course_preamble.tex}
\title[Policy Function Iteration]{Lecture: Policy Function Iteration (PFI)}
\newcommand{\E}{\mathbb E}
\newcommand{\Sset}{\mathcal S}
\newcommand{\Aset}{\mathcal A}
\tikzset{valueState/.style={draw=SessionTeal,fill=SessionTealLight,rounded corners=1mm,minimum width=15mm,minimum height=9mm,align=center},valueFlow/.style={-{Stealth[length=2mm]},draw=SessionCharcoal,line width=.8pt},stateCircle/.style={circle,draw=SessionTeal,fill=white,minimum size=6mm,inner sep=1pt},actionCircle/.style={circle,draw=SessionTeal,fill=SessionTeal,minimum size=3mm,inner sep=0pt}}
% The maze geometry and reward convention match the Week 1 example.
\definecolor{ValueMazeWall}{HTML}{928580}
\definecolor{MazeStartFill}{HTML}{F5DEA2}
\definecolor{MazeGoalFill}{HTML}{B5DDD5}
\newcommand{\ValueMazeBase}{%
  \fill[white] (0,0) rectangle (9,6);
  \fill[MazeStartFill] (0,3) rectangle +(1,1);
  \fill[MazeGoalFill] (8,5) rectangle +(1,1);
  \foreach \x/\y in {2/2,2/3,2/4,7/3,7/4,7/5,5/1}
    \fill[ValueMazeWall] (\x,\y) rectangle +(1,1);
  \draw[SessionLine,line width=.3pt,step=1] (0,0) grid (9,6);
  \draw[SessionCharcoal,line width=.7pt] (0,0) rectangle (9,6);
  \node[anchor=north west,inner sep=0pt,font=\sffamily\fontsize{6}{6}\selectfont,text=SessionMuted!65] at (0.09,5.92) {1};
  \node[anchor=north west,inner sep=0pt,font=\sffamily\fontsize{6}{6}\selectfont,text=SessionMuted!65] at (1.09,5.92) {2};
  \node[anchor=north west,inner sep=0pt,font=\sffamily\fontsize{6}{6}\selectfont,text=SessionMuted!65] at (2.09,5.92) {3};
  \node[anchor=north west,inner sep=0pt,font=\sffamily\fontsize{6}{6}\selectfont,text=SessionMuted!65] at (3.09,5.92) {4};
  \node[anchor=north west,inner sep=0pt,font=\sffamily\fontsize{6}{6}\selectfont,text=SessionMuted!65] at (4.09,5.92) {5};
  \node[anchor=north west,inner sep=0pt,font=\sffamily\fontsize{6}{6}\selectfont,text=SessionMuted!65] at (5.09,5.92) {6};
  \node[anchor=north west,inner sep=0pt,font=\sffamily\fontsize{6}{6}\selectfont,text=SessionMuted!65] at (6.09,5.92) {7};
  \node[anchor=north west,inner sep=0pt,font=\sffamily\fontsize{6}{6}\selectfont,text=SessionMuted!65] at (8.09,5.92) {8};
  \node[anchor=north west,inner sep=0pt,font=\sffamily\fontsize{6}{6}\selectfont,text=SessionMuted!65] at (0.09,4.92) {9};
  \node[anchor=north west,inner sep=0pt,font=\sffamily\fontsize{6}{6}\selectfont,text=SessionMuted!65] at (1.09,4.92) {10};
  \node[anchor=north west,inner sep=0pt,font=\sffamily\fontsize{6}{6}\selectfont,text=SessionMuted!65] at (3.09,4.92) {11};
  \node[anchor=north west,inner sep=0pt,font=\sffamily\fontsize{6}{6}\selectfont,text=SessionMuted!65] at (4.09,4.92) {12};
  \node[anchor=north west,inner sep=0pt,font=\sffamily\fontsize{6}{6}\selectfont,text=SessionMuted!65] at (5.09,4.92) {13};
  \node[anchor=north west,inner sep=0pt,font=\sffamily\fontsize{6}{6}\selectfont,text=SessionMuted!65] at (6.09,4.92) {14};
  \node[anchor=north west,inner sep=0pt,font=\sffamily\fontsize{6}{6}\selectfont,text=SessionMuted!65] at (8.09,4.92) {15};
  \node[anchor=north west,inner sep=0pt,font=\sffamily\fontsize{6}{6}\selectfont,text=SessionMuted!65] at (0.09,3.92) {16};
  \node[anchor=north west,inner sep=0pt,font=\sffamily\fontsize{6}{6}\selectfont,text=SessionMuted!65] at (1.09,3.92) {17};
  \node[anchor=north west,inner sep=0pt,font=\sffamily\fontsize{6}{6}\selectfont,text=SessionMuted!65] at (3.09,3.92) {18};
  \node[anchor=north west,inner sep=0pt,font=\sffamily\fontsize{6}{6}\selectfont,text=SessionMuted!65] at (4.09,3.92) {19};
  \node[anchor=north west,inner sep=0pt,font=\sffamily\fontsize{6}{6}\selectfont,text=SessionMuted!65] at (5.09,3.92) {20};
  \node[anchor=north west,inner sep=0pt,font=\sffamily\fontsize{6}{6}\selectfont,text=SessionMuted!65] at (6.09,3.92) {21};
  \node[anchor=north west,inner sep=0pt,font=\sffamily\fontsize{6}{6}\selectfont,text=SessionMuted!65] at (8.09,3.92) {22};
  \node[anchor=north west,inner sep=0pt,font=\sffamily\fontsize{6}{6}\selectfont,text=SessionMuted!65] at (0.09,2.92) {23};
  \node[anchor=north west,inner sep=0pt,font=\sffamily\fontsize{6}{6}\selectfont,text=SessionMuted!65] at (1.09,2.92) {24};
  \node[anchor=north west,inner sep=0pt,font=\sffamily\fontsize{6}{6}\selectfont,text=SessionMuted!65] at (3.09,2.92) {25};
  \node[anchor=north west,inner sep=0pt,font=\sffamily\fontsize{6}{6}\selectfont,text=SessionMuted!65] at (4.09,2.92) {26};
  \node[anchor=north west,inner sep=0pt,font=\sffamily\fontsize{6}{6}\selectfont,text=SessionMuted!65] at (5.09,2.92) {27};
  \node[anchor=north west,inner sep=0pt,font=\sffamily\fontsize{6}{6}\selectfont,text=SessionMuted!65] at (6.09,2.92) {28};
  \node[anchor=north west,inner sep=0pt,font=\sffamily\fontsize{6}{6}\selectfont,text=SessionMuted!65] at (7.09,2.92) {29};
  \node[anchor=north west,inner sep=0pt,font=\sffamily\fontsize{6}{6}\selectfont,text=SessionMuted!65] at (8.09,2.92) {30};
  \node[anchor=north west,inner sep=0pt,font=\sffamily\fontsize{6}{6}\selectfont,text=SessionMuted!65] at (0.09,1.92) {31};
  \node[anchor=north west,inner sep=0pt,font=\sffamily\fontsize{6}{6}\selectfont,text=SessionMuted!65] at (1.09,1.92) {32};
  \node[anchor=north west,inner sep=0pt,font=\sffamily\fontsize{6}{6}\selectfont,text=SessionMuted!65] at (2.09,1.92) {33};
  \node[anchor=north west,inner sep=0pt,font=\sffamily\fontsize{6}{6}\selectfont,text=SessionMuted!65] at (3.09,1.92) {34};
  \node[anchor=north west,inner sep=0pt,font=\sffamily\fontsize{6}{6}\selectfont,text=SessionMuted!65] at (4.09,1.92) {35};
  \node[anchor=north west,inner sep=0pt,font=\sffamily\fontsize{6}{6}\selectfont,text=SessionMuted!65] at (6.09,1.92) {36};
  \node[anchor=north west,inner sep=0pt,font=\sffamily\fontsize{6}{6}\selectfont,text=SessionMuted!65] at (7.09,1.92) {37};
  \node[anchor=north west,inner sep=0pt,font=\sffamily\fontsize{6}{6}\selectfont,text=SessionMuted!65] at (8.09,1.92) {38};
  \node[anchor=north west,inner sep=0pt,font=\sffamily\fontsize{6}{6}\selectfont,text=SessionMuted!65] at (0.09,0.92) {39};
  \node[anchor=north west,inner sep=0pt,font=\sffamily\fontsize{6}{6}\selectfont,text=SessionMuted!65] at (1.09,0.92) {40};
  \node[anchor=north west,inner sep=0pt,font=\sffamily\fontsize{6}{6}\selectfont,text=SessionMuted!65] at (2.09,0.92) {41};
  \node[anchor=north west,inner sep=0pt,font=\sffamily\fontsize{6}{6}\selectfont,text=SessionMuted!65] at (3.09,0.92) {42};
  \node[anchor=north west,inner sep=0pt,font=\sffamily\fontsize{6}{6}\selectfont,text=SessionMuted!65] at (4.09,0.92) {43};
  \node[anchor=north west,inner sep=0pt,font=\sffamily\fontsize{6}{6}\selectfont,text=SessionMuted!65] at (5.09,0.92) {44};
  \node[anchor=north west,inner sep=0pt,font=\sffamily\fontsize{6}{6}\selectfont,text=SessionMuted!65] at (6.09,0.92) {45};
  \node[anchor=north west,inner sep=0pt,font=\sffamily\fontsize{6}{6}\selectfont,text=SessionMuted!65] at (7.09,0.92) {46};
  \node[anchor=north west,inner sep=0pt,font=\sffamily\fontsize{6}{6}\selectfont,text=SessionMuted!65] at (8.09,0.92) {47};
}
\newcommand{\MazeLegend}{%
\par\smallskip{\scriptsize
\tikz[baseline=-.5ex]\draw[SessionLine,fill=MazeStartFill] (0,0) rectangle (.22,.16);\ Start (16)\qquad
\tikz[baseline=-.5ex]\draw[SessionLine,fill=MazeGoalFill] (0,0) rectangle (.22,.16);\ Goal (8)}
}

% Centered arrows overlay cell text, matching the continuation diagram.
\newcommand{\MazePolicyPath}{%
  \draw[valueFlow,draw=CourseVariableColor,shorten <=2pt,shorten >=2pt] (0.5,3.5) -- (0.5,2.5);
  \draw[valueFlow,draw=CourseVariableColor,shorten <=2pt,shorten >=2pt] (0.5,2.5) -- (0.5,1.5);
  \draw[valueFlow,draw=CourseVariableColor,shorten <=2pt,shorten >=2pt] (0.5,1.5) -- (1.5,1.5);
  \draw[valueFlow,draw=CourseVariableColor,shorten <=2pt,shorten >=2pt] (1.5,1.5) -- (2.5,1.5);
  \draw[valueFlow,draw=CourseVariableColor,shorten <=2pt,shorten >=2pt] (2.5,1.5) -- (3.5,1.5);
  \draw[valueFlow,draw=CourseVariableColor,shorten <=2pt,shorten >=2pt] (3.5,1.5) -- (3.5,2.5);
  \draw[valueFlow,draw=CourseVariableColor,shorten <=2pt,shorten >=2pt] (3.5,2.5) -- (3.5,3.5);
  \draw[valueFlow,draw=CourseVariableColor,shorten <=2pt,shorten >=2pt] (3.5,3.5) -- (4.5,3.5);
  \draw[valueFlow,draw=CourseVariableColor,shorten <=2pt,shorten >=2pt] (4.5,3.5) -- (5.5,3.5);
  \draw[valueFlow,draw=CourseVariableColor,shorten <=2pt,shorten >=2pt] (5.5,3.5) -- (6.5,3.5);
  \draw[valueFlow,draw=CourseVariableColor,shorten <=2pt,shorten >=2pt] (6.5,3.5) -- (6.5,2.5);
  \draw[valueFlow,draw=CourseVariableColor,shorten <=2pt,shorten >=2pt] (6.5,2.5) -- (7.5,2.5);
  \draw[valueFlow,draw=CourseVariableColor,shorten <=2pt,shorten >=2pt] (7.5,2.5) -- (8.5,2.5);
  \draw[valueFlow,draw=CourseVariableColor,shorten <=2pt,shorten >=2pt] (8.5,2.5) -- (8.5,3.5);
  \draw[valueFlow,draw=CourseVariableColor,shorten <=2pt,shorten >=2pt] (8.5,3.5) -- (8.5,4.5);
  \draw[valueFlow,draw=CourseVariableColor,shorten <=2pt,shorten >=2pt] (8.5,4.5) -- (8.5,5.5);
}

\newcommand{\MazeDiscountedValues}{%
\node[font=\sffamily\footnotesize,text=SessionMuted!55,inner sep=0pt] at (0.5,5.5) {-7.7};
\node[font=\sffamily\footnotesize,text=SessionMuted!55,inner sep=0pt] at (1.5,5.5) {-7.5};
\node[font=\sffamily\footnotesize,text=SessionMuted!55,inner sep=0pt] at (2.5,5.5) {-7.2};
\node[font=\sffamily\footnotesize,text=SessionMuted!55,inner sep=0pt] at (3.5,5.5) {-6.9};
\node[font=\sffamily\footnotesize,text=SessionMuted!55,inner sep=0pt] at (4.5,5.5) {-6.5};
\node[font=\sffamily\footnotesize,text=SessionMuted!55,inner sep=0pt] at (5.5,5.5) {-6.1};
\node[font=\sffamily\footnotesize,text=SessionMuted!55,inner sep=0pt] at (6.5,5.5) {-5.7};
\node[font=\sffamily\footnotesize,text=SessionMuted!55,inner sep=0pt] at (8.5,5.5) {0.0};
\node[font=\sffamily\footnotesize,text=SessionMuted!55,inner sep=0pt] at (0.5,4.5) {-7.9};
\node[font=\sffamily\footnotesize,text=SessionMuted!55,inner sep=0pt] at (1.5,4.5) {-7.7};
\node[font=\sffamily\footnotesize,text=SessionMuted!55,inner sep=0pt] at (3.5,4.5) {-6.5};
\node[font=\sffamily\footnotesize,text=SessionMuted!55,inner sep=0pt] at (4.5,4.5) {-6.1};
\node[font=\sffamily\footnotesize,text=SessionMuted!55,inner sep=0pt] at (5.5,4.5) {-5.7};
\node[font=\sffamily\footnotesize,text=SessionMuted!55,inner sep=0pt] at (6.5,4.5) {-5.2};
\node[font=\sffamily\footnotesize,text=SessionMuted!55,inner sep=0pt] at (8.5,4.5) {-1.0};
\node[font=\sffamily\footnotesize,text=SessionMuted!55,inner sep=0pt] at (0.5,3.5) {-8.1};
\node[font=\sffamily\footnotesize,text=SessionMuted!55,inner sep=0pt] at (1.5,3.5) {-7.9};
\node[font=\sffamily\footnotesize,text=SessionMuted!55,inner sep=0pt] at (3.5,3.5) {-6.1};
\node[font=\sffamily\footnotesize,text=SessionMuted!55,inner sep=0pt] at (4.5,3.5) {-5.7};
\node[font=\sffamily\footnotesize,text=SessionMuted!55,inner sep=0pt] at (5.5,3.5) {-5.2};
\node[font=\sffamily\footnotesize,text=SessionMuted!55,inner sep=0pt] at (6.5,3.5) {-4.7};
\node[font=\sffamily\footnotesize,text=SessionMuted!55,inner sep=0pt] at (8.5,3.5) {-1.9};
\node[font=\sffamily\footnotesize,text=SessionMuted!55,inner sep=0pt] at (0.5,2.5) {-7.9};
\node[font=\sffamily\footnotesize,text=SessionMuted!55,inner sep=0pt] at (1.5,2.5) {-7.7};
\node[font=\sffamily\footnotesize,text=SessionMuted!55,inner sep=0pt] at (3.5,2.5) {-6.5};
\node[font=\sffamily\footnotesize,text=SessionMuted!55,inner sep=0pt] at (4.5,2.5) {-5.2};
\node[font=\sffamily\footnotesize,text=SessionMuted!55,inner sep=0pt] at (5.5,2.5) {-4.7};
\node[font=\sffamily\footnotesize,text=SessionMuted!55,inner sep=0pt] at (6.5,2.5) {-4.1};
\node[font=\sffamily\footnotesize,text=SessionMuted!55,inner sep=0pt] at (7.5,2.5) {-3.4};
\node[font=\sffamily\footnotesize,text=SessionMuted!55,inner sep=0pt] at (8.5,2.5) {-2.7};
\node[font=\sffamily\footnotesize,text=SessionMuted!55,inner sep=0pt] at (0.5,1.5) {-7.7};
\node[font=\sffamily\footnotesize,text=SessionMuted!55,inner sep=0pt] at (1.5,1.5) {-7.5};
\node[font=\sffamily\footnotesize,text=SessionMuted!55,inner sep=0pt] at (2.5,1.5) {-7.2};
\node[font=\sffamily\footnotesize,text=SessionMuted!55,inner sep=0pt] at (3.5,1.5) {-6.9};
\node[font=\sffamily\footnotesize,text=SessionMuted!55,inner sep=0pt] at (4.5,1.5) {-5.7};
\node[font=\sffamily\footnotesize,text=SessionMuted!55,inner sep=0pt] at (6.5,1.5) {-4.7};
\node[font=\sffamily\footnotesize,text=SessionMuted!55,inner sep=0pt] at (7.5,1.5) {-4.1};
\node[font=\sffamily\footnotesize,text=SessionMuted!55,inner sep=0pt] at (8.5,1.5) {-3.4};
\node[font=\sffamily\footnotesize,text=SessionMuted!55,inner sep=0pt] at (0.5,0.5) {-7.9};
\node[font=\sffamily\footnotesize,text=SessionMuted!55,inner sep=0pt] at (1.5,0.5) {-7.7};
\node[font=\sffamily\footnotesize,text=SessionMuted!55,inner sep=0pt] at (2.5,0.5) {-7.5};
\node[font=\sffamily\footnotesize,text=SessionMuted!55,inner sep=0pt] at (3.5,0.5) {-7.2};
\node[font=\sffamily\footnotesize,text=SessionMuted!55,inner sep=0pt] at (4.5,0.5) {-6.1};
\node[font=\sffamily\footnotesize,text=SessionMuted!55,inner sep=0pt] at (5.5,0.5) {-5.7};
\node[font=\sffamily\footnotesize,text=SessionMuted!55,inner sep=0pt] at (6.5,0.5) {-5.2};
\node[font=\sffamily\footnotesize,text=SessionMuted!55,inner sep=0pt] at (7.5,0.5) {-4.7};
\node[font=\sffamily\footnotesize,text=SessionMuted!55,inner sep=0pt] at (8.5,0.5) {-4.1};
}
\newcommand{\MazeImprovedValues}{%
\node[font=\sffamily\footnotesize,text=SessionMuted!55,inner sep=0pt] at (0.5,5.5) {-7.7};
\node[font=\sffamily\footnotesize,text=SessionMuted!55,inner sep=0pt] at (1.5,5.5) {-7.5};
\node[font=\sffamily\footnotesize,text=SessionMuted!55,inner sep=0pt] at (2.5,5.5) {-7.2};
\node[font=\sffamily\footnotesize,text=SessionMuted!55,inner sep=0pt] at (3.5,5.5) {-6.9};
\node[font=\sffamily\footnotesize,text=SessionMuted!55,inner sep=0pt] at (4.5,5.5) {-6.5};
\node[font=\sffamily\footnotesize,text=SessionMuted!55,inner sep=0pt] at (5.5,5.5) {-6.1};
\node[font=\sffamily\footnotesize,text=SessionMuted!55,inner sep=0pt] at (6.5,5.5) {-5.7};
\node[font=\sffamily\footnotesize,text=SessionMuted!55,inner sep=0pt] at (8.5,5.5) {0.0};
\node[font=\sffamily\footnotesize,text=SessionMuted!55,inner sep=0pt] at (0.5,4.5) {-7.9};
\node[font=\sffamily\footnotesize,text=SessionMuted!55,inner sep=0pt] at (1.5,4.5) {-7.7};
\node[font=\sffamily\footnotesize,text=SessionMuted!55,inner sep=0pt] at (3.5,4.5) {-6.5};
\node[font=\sffamily\footnotesize,text=SessionMuted!55,inner sep=0pt] at (4.5,4.5) {-6.1};
\node[font=\sffamily\footnotesize,text=SessionMuted!55,inner sep=0pt] at (5.5,4.5) {-5.7};
\node[font=\sffamily\footnotesize,text=SessionMuted!55,inner sep=0pt] at (6.5,4.5) {-5.2};
\node[font=\sffamily\footnotesize,text=SessionMuted!55,inner sep=0pt] at (8.5,4.5) {-1.0};
\node[font=\sffamily\footnotesize,text=SessionMuted!55,inner sep=0pt] at (0.5,3.5) {-7.7};
\node[font=\sffamily\footnotesize,text=SessionMuted!55,inner sep=0pt] at (1.5,3.5) {-7.5};
\node[font=\sffamily\footnotesize,text=SessionMuted!55,inner sep=0pt] at (3.5,3.5) {-6.1};
\node[font=\sffamily\footnotesize,text=SessionMuted!55,inner sep=0pt] at (4.5,3.5) {-5.7};
\node[font=\sffamily\footnotesize,text=SessionMuted!55,inner sep=0pt] at (5.5,3.5) {-5.2};
\node[font=\sffamily\footnotesize,text=SessionMuted!55,inner sep=0pt] at (6.5,3.5) {-4.7};
\node[font=\sffamily\footnotesize,text=SessionMuted!55,inner sep=0pt] at (8.5,3.5) {-1.9};
\node[font=\sffamily\footnotesize,text=SessionMuted!55,inner sep=0pt] at (0.5,2.5) {-7.5};
\node[font=\sffamily\footnotesize,text=SessionMuted!55,inner sep=0pt] at (1.5,2.5) {-7.2};
\node[font=\sffamily\footnotesize,text=SessionMuted!55,inner sep=0pt] at (3.5,2.5) {-5.7};
\node[font=\sffamily\footnotesize,text=SessionMuted!55,inner sep=0pt] at (4.5,2.5) {-5.2};
\node[font=\sffamily\footnotesize,text=SessionMuted!55,inner sep=0pt] at (5.5,2.5) {-4.7};
\node[font=\sffamily\footnotesize,text=SessionMuted!55,inner sep=0pt] at (6.5,2.5) {-4.1};
\node[font=\sffamily\footnotesize,text=SessionMuted!55,inner sep=0pt] at (7.5,2.5) {-3.4};
\node[font=\sffamily\footnotesize,text=SessionMuted!55,inner sep=0pt] at (8.5,2.5) {-2.7};
\node[font=\sffamily\footnotesize,text=SessionMuted!55,inner sep=0pt] at (0.5,1.5) {-7.2};
\node[font=\sffamily\footnotesize,text=SessionMuted!55,inner sep=0pt] at (1.5,1.5) {-6.9};
\node[font=\sffamily\footnotesize,text=SessionMuted!55,inner sep=0pt] at (2.5,1.5) {-6.5};
\node[font=\sffamily\footnotesize,text=SessionMuted!55,inner sep=0pt] at (3.5,1.5) {-6.1};
\node[font=\sffamily\footnotesize,text=SessionMuted!55,inner sep=0pt] at (4.5,1.5) {-5.7};
\node[font=\sffamily\footnotesize,text=SessionMuted!55,inner sep=0pt] at (6.5,1.5) {-4.7};
\node[font=\sffamily\footnotesize,text=SessionMuted!55,inner sep=0pt] at (7.5,1.5) {-4.1};
\node[font=\sffamily\footnotesize,text=SessionMuted!55,inner sep=0pt] at (8.5,1.5) {-3.4};
\node[font=\sffamily\footnotesize,text=SessionMuted!55,inner sep=0pt] at (0.5,0.5) {-7.5};
\node[font=\sffamily\footnotesize,text=SessionMuted!55,inner sep=0pt] at (1.5,0.5) {-7.2};
\node[font=\sffamily\footnotesize,text=SessionMuted!55,inner sep=0pt] at (2.5,0.5) {-6.9};
\node[font=\sffamily\footnotesize,text=SessionMuted!55,inner sep=0pt] at (3.5,0.5) {-6.5};
\node[font=\sffamily\footnotesize,text=SessionMuted!55,inner sep=0pt] at (4.5,0.5) {-6.1};
\node[font=\sffamily\footnotesize,text=SessionMuted!55,inner sep=0pt] at (5.5,0.5) {-5.7};
\node[font=\sffamily\footnotesize,text=SessionMuted!55,inner sep=0pt] at (6.5,0.5) {-5.2};
\node[font=\sffamily\footnotesize,text=SessionMuted!55,inner sep=0pt] at (7.5,0.5) {-4.7};
\node[font=\sffamily\footnotesize,text=SessionMuted!55,inner sep=0pt] at (8.5,0.5) {-4.1};
}
\newcommand{\MazeImprovedPath}{%
\draw[valueFlow,draw=CourseVariableColor,shorten <=2pt,shorten >=2pt] (0.5,3.5) -- (0.5,2.5);
\draw[valueFlow,draw=CourseVariableColor,shorten <=2pt,shorten >=2pt] (0.5,2.5) -- (0.5,1.5);
\draw[valueFlow,draw=CourseVariableColor,shorten <=2pt,shorten >=2pt] (0.5,1.5) -- (1.5,1.5);
\draw[valueFlow,draw=CourseVariableColor,shorten <=2pt,shorten >=2pt] (1.5,1.5) -- (2.5,1.5);
\draw[valueFlow,draw=CourseVariableColor,shorten <=2pt,shorten >=2pt] (2.5,1.5) -- (3.5,1.5);
\draw[valueFlow,draw=CourseVariableColor,shorten <=2pt,shorten >=2pt] (3.5,1.5) -- (3.5,2.5);
\draw[valueFlow,draw=SessionRed,shorten <=2pt,shorten >=2pt] (3.5,2.5) -- (4.5,2.5);
\draw[valueFlow,draw=SessionRed,shorten <=2pt,shorten >=2pt] (4.5,2.5) -- (5.5,2.5);
\draw[valueFlow,draw=SessionRed,shorten <=2pt,shorten >=2pt] (5.5,2.5) -- (6.5,2.5);
\draw[valueFlow,draw=CourseVariableColor,shorten <=2pt,shorten >=2pt] (6.5,2.5) -- (7.5,2.5);
\draw[valueFlow,draw=CourseVariableColor,shorten <=2pt,shorten >=2pt] (7.5,2.5) -- (8.5,2.5);
\draw[valueFlow,draw=CourseVariableColor,shorten <=2pt,shorten >=2pt] (8.5,2.5) -- (8.5,3.5);
\draw[valueFlow,draw=CourseVariableColor,shorten <=2pt,shorten >=2pt] (8.5,3.5) -- (8.5,4.5);
\draw[valueFlow,draw=CourseVariableColor,shorten <=2pt,shorten >=2pt] (8.5,4.5) -- (8.5,5.5);
}

\begin{document}
\makecoursetitle

\begin{frame}{Lecture Overview}
\begin{itemize}
\item \textbf{Policy Evaluation}
  \begin{itemize}
  \item The policy evaluation operator
  \item Repeated updates and a linear system
  \end{itemize}
\item \textbf{Policy Function Iteration}
  \begin{itemize}
  \item Evaluation and improvement in the maze
  \item The algorithm and its optimality criterion
  \end{itemize}
\item \textbf{Computational Choices}
  \begin{itemize}
  \item Generalized and modified policy iteration
  \item Computational cost and continuous states
  \end{itemize}
\end{itemize}
\end{frame}

\begin{frame}{From Value Iteration to Policy Iteration}
\small
Given the current value estimate \(\modelvar{V_k}\), a VFI update separates into two steps:
\begin{enumerate}
\item \term{Policy improvement}: choose a policy greedy with respect to \(\modelvar{V_k}\):
\[
\sigma_{k+1}(s)\in\argmax_{a\in\Aset(s)}
\{r(s,a)+\beta\E[V_k(s')\mid s,a]\}.
\]
\item \term{Policy evaluation update}: hold \(\modelvar{\sigma_{k+1}}\) fixed and update the value once:
\[
V_{k+1}(s)=r(s,\sigma_{k+1}(s))
+\beta\E[V_k(s')\mid s,a=\sigma_{k+1}(s)].
\]
\end{enumerate}
Because the selected action is a maximizer, these steps give exactly the VFI update.
\par\medskip
\term{Policy evaluation} computes a fixed policy's value. VFI takes only one update; generally, \(\modelvar{V_{k+1}\ne V^{\sigma_{k+1}}}\).
\par\medskip
\term{Policy function iteration (PFI)} instead completes policy evaluation to obtain \(\modelvar{V^{\sigma_{k+1}}}\) before the next policy improvement.
\end{frame}

\begin{frame}{The Policy Evaluation Operator}
\small
For a fixed policy \(\modelvar{\sigma}\), the \term{policy evaluation operator} \(\modelvar{T_\sigma}\) maps a candidate continuation-value function \(\modelvar{V}\) to an updated function:
\[
(T_\sigma V)(s)=r(s,\sigma(s))+\beta\E[V(s')\mid s,a=\sigma(s)].
\]
\begin{itemize}
\item The action is supplied by \(\modelvar{\sigma(s)}\); this update does not choose a new action.
\item The input \(\modelvar{V}\) may be any bounded candidate function. It need not equal \(\modelvar{V^\sigma}\).
\end{itemize}
\par\medskip
\term{Policy evaluation} means finding the fixed point:
\[V^\sigma=T_\sigma V^\sigma.\]
One application \(\modelvar{T_\sigma V}\) is one update; repeated updates compute the policy's value.
\end{frame}

\begin{frame}{Iterative Policy Evaluation}
\small
Hold \(\modelvar{\sigma}\) fixed. Choose \(\modelvar{V_0}\), and update all states using the same previous function:
\[
V_{j+1}(s)=(T_\sigma V_j)(s)
=r(s,\sigma(s))+\beta\sum_{s'\in\Sset}P(s'\mid s,\sigma(s))V_j(s').
\]
The index \(\modelvar{j}\) counts evaluation updates, not changes of policy.
\par\medskip
As with the Bellman optimality operator, discounting gives a contraction:
\[\|T_\sigma V_1-T_\sigma V_2\|_\infty\leq\beta\|V_1-V_2\|_\infty.\]
Hence \(\modelvar{V_j\to V^\sigma}\). For a target evaluation error \(\modelvar{\varepsilon_{\rm eval}>0}\), compute
\[
d_j=\|V_{j+1}-V_j\|_\infty,
\qquad
\|V_j-V^\sigma\|_\infty\leq\frac{d_j}{1-\beta}.
\]
Stop and return \(\modelvar{V_j}\) when \(\modelvar{d_j\leq(1-\beta)\varepsilon_{\rm eval}}\).
\end{frame}

\begin{frame}{Policy Evaluation as a Linear System}
\small
For finitely many states, collect the rewards and transition probabilities selected by \(\modelvar{\sigma}\):
\[
r^\sigma(s)=r(s,\sigma(s)),\qquad
P^\sigma(s,s')=P(s'\mid s,\sigma(s)).
\]
Each row of \(\modelvar{P^\sigma}\) describes the next-state probabilities under the policy's chosen action.
\par\medskip
The policy's Bellman equation becomes
\[
V^\sigma=r^\sigma+\beta P^\sigma V^\sigma
\qquad\Longleftrightarrow\qquad
(I-\beta P^\sigma)V^\sigma=r^\sigma.
\]
\begin{itemize}
\item \(\modelvar{I}\) is the identity matrix. There is one equation and one unknown value per state.
\item Solve this linear system directly, or use repeated evaluation updates.
\item A direct solve avoids an inner iteration; its cost grows with the number of states.
\end{itemize}
\end{frame}

\begin{frame}{Example: Evaluating the Maze Policy}
\small
Use the same policy \(\modelvar{\sigma}\) as before, now with \(\modelvar{\beta=0.9}\).
\begin{columns}[c]
\begin{column}{.58\textwidth}
\begin{figure}\centering
\begin{tikzpicture}[x=.71cm,y=.71cm]
\ValueMazeBase\MazeDiscountedValues\MazePolicyPath
\end{tikzpicture}\MazeLegend
\caption{Discounted values under \(\sigma\).}
\end{figure}
\end{column}
\begin{column}{.40\textwidth}
\begin{itemize}
\item Each move gives \(\modelvar{-1}\), including entry into cell 8; rewards afterward are zero.
\item At cell 25, the policy chooses U to cell 18.
\end{itemize}
\[
\begin{aligned}
V^\sigma(8)&=0,\\
V^\sigma(25)&=-1+0.9V^\sigma(18)\\
&\approx -6.513.
\end{aligned}
\]
\par\smallskip
At the starting cell, \(\modelvar{V^\sigma(16)\approx-8.147}\).
\par\smallskip
\footnotesize Cell values are rounded to one decimal place.
\end{column}
\end{columns}
\end{frame}

\begin{frame}{Example: Improving and Re-evaluating the Policy}
\small
Change only cell 25 to R: \(\modelvar{\widetilde\sigma(25)=\mathrm R}\); keep \(\modelvar{\widetilde\sigma(s)=\sigma(s)}\) elsewhere.
\begin{columns}[c]
\begin{column}{.58\textwidth}
\begin{figure}\centering
\begin{tikzpicture}[x=.71cm,y=.71cm]
\ValueMazeBase
\fill[CourseVariableColor!20] (3.01,2.01) rectangle (3.99,2.99);
\node[anchor=north west,inner sep=0pt,font=\sffamily\bfseries\fontsize{6}{6}\selectfont,text=CourseVariableColor] at (3.09,2.92) {25};
\MazeImprovedValues\MazeImprovedPath
\end{tikzpicture}\MazeLegend
\caption{Re-evaluated values under \(\widetilde\sigma\).}
\end{figure}
\end{column}
\begin{column}{.40\textwidth}
\textbf{Policy improvement using \(\modelvar{V^\sigma}\):}
\[
\begin{aligned}
Q^\sigma(25,\mathrm U)&\approx-6.513,\\
Q^\sigma(25,\mathrm R)&\approx-5.695.
\end{aligned}
\]
\textbf{Policy evaluation for the new policy:}
\[
\begin{aligned}
V^{\widetilde\sigma}(25)&\approx-5.695,\\
V^{\widetilde\sigma}(16)&\approx-7.712.
\end{aligned}
\]
The shortcut also improves values at states whose paths pass through cell 25.
\end{column}
\end{columns}
\end{frame}

\begin{frame}{Policy Function Iteration (PFI) Algorithm}
\small
\begin{enumerate}
\setcounter{enumi}{-1}
\item \textbf{Input.} Finite state set \(\modelvar{\Sset}\), finite nonempty feasible action sets \(\modelvar{\Aset(s)}\), rewards \(\modelvar{r(s,a)}\), transition probabilities \(\modelvar{P(s'\mid s,a)}\), and discount factor \(\modelvar{\beta}\).
\item \textbf{Initialize.} Choose a feasible policy \(\modelvar{\sigma_0(s)\in\Aset(s)}\) for every state. Set \(\modelvar{k=0}\).
\item \textbf{Policy evaluation.} Form \(\modelvar{r^{\sigma_k}}\) and \(\modelvar{P^{\sigma_k}}\), then solve
\[(I-\beta P^{\sigma_k})V^{\sigma_k}=r^{\sigma_k}.\]
\item \textbf{Compute action values.} For every feasible \(\modelvar{(s,a)}\), compute
\[Q^{\sigma_k}(s,a)=r(s,a)+\beta\sum_{s'\in\Sset}P(s'\mid s,a)V^{\sigma_k}(s').\]
Use this same evaluated \(\modelvar{V^{\sigma_k}}\) for every action comparison.
\end{enumerate}
\end{frame}

\begin{frame}{PFI Algorithm (continued)}
\small
\begin{enumerate}
\setcounter{enumi}{3}
\item \textbf{Policy improvement.} For every state, select
\[\sigma_{k+1}(s)\in\argmax_{a\in\Aset(s)}Q^{\sigma_k}(s,a).\]
Keep \(\modelvar{\sigma_k(s)}\) if it is already a maximizer; otherwise use a fixed action order to break ties.
\item \textbf{Check the policy.}
\begin{itemize}
\item If \(\modelvar{\sigma_{k+1}(s)=\sigma_k(s)}\) at every state, continue to \textbf{Step 6}.
\item Otherwise, set \(\modelvar{k\leftarrow k+1}\) and return to \textbf{Step 2}.
\end{itemize}
\item \textbf{Return and stop.} Output the policy and its evaluated value:
\[\sigma^*=\sigma_k,\qquad V^*=V^{\sigma_k}.\]
\end{enumerate}
\par\smallskip
\footnotesize The exact optimality statement uses exact policy evaluation. A numerical solve introduces rounding error; truncated evaluation introduces an additional approximation.
\end{frame}

\begin{frame}{Policy Improvement and Optimality}
\small
Every greedy improvement satisfies the policy improvement theorem:
\[
Q^{\sigma_k}(s,\sigma_{k+1}(s))\geq V^{\sigma_k}(s)\quad\forall s
\quad\Longrightarrow\quad
V^{\sigma_{k+1}}(s)\geq V^{\sigma_k}(s)\quad\forall s.
\]
\par\smallskip
When the policy does not change, it is greedy with respect to its own value:
\[
V^{\sigma_k}(s)=\max_{a\in\Aset(s)}Q^{\sigma_k}(s,a)
=(TV^{\sigma_k})(s).
\]
Thus \(\modelvar{V^{\sigma_k}=V^*}\), by uniqueness of the Bellman optimality fixed point.
\par\medskip
\begin{itemize}
\item Retaining the current action in a tie prevents changes between equally good actions.
\item If the policy changes, its value strictly increases at some state.
\item With finitely many deterministic policies, exact PFI therefore terminates at an optimal policy after finitely many improvements.
\end{itemize}
\end{frame}

\begin{frame}{Generalized Policy Iteration}
\small
\term{Generalized policy iteration (GPI)} lets policy evaluation and policy improvement interact, allowing evaluation to stop before convergence.
\begin{figure}
\centering
\begin{tikzpicture}[x=.93cm,y=.93cm,
  gpiEval/.style={-{Stealth[length=1.5mm]},CourseVariableColor,line width=.8pt},
  gpiImprove/.style={-{Stealth[length=1.5mm]},SessionTeal,line width=.8pt}]
\begin{scope}[xshift=0cm]
\node[font=\footnotesize\bfseries] at (2.45,2.65) {Full policy evaluation};
\draw[SessionMuted,line width=.6pt] (0,2.1) -- (4.9,1.05) -- (0,0);
\node[font=\scriptsize,anchor=south,fill=SessionPaper,inner sep=1pt] at (2.1,1.77) {\(V=V^\sigma\)};
\node[font=\scriptsize,anchor=north] at (2.25,.34) {\(\sigma\) greedy w.r.t. \(V\)};
\fill[SessionCharcoal] (4.9,1.05) circle (1.1pt);
\node[font=\scriptsize,anchor=west] at (4.96,1.05) {\(V^*,\sigma^*\)};
\fill[SessionCharcoal] (0.15,0.65) circle (1pt);
\draw[gpiEval] (0.1500,0.6500) -- (0.8625,1.9152);
\draw[gpiImprove] (0.8625,1.9152) -- (2.0738,0.4444);
\draw[gpiEval] (2.0738,0.4444) -- (2.4977,1.5648);
\draw[gpiImprove] (2.4977,1.5648) -- (3.2184,0.6897);
\draw[gpiEval] (3.2184,0.6897) -- (3.4706,1.3563);
\draw[gpiImprove] (3.4706,1.3563) -- (3.8994,0.8356);
\draw[gpiEval] (3.8994,0.8356) -- (4.0495,1.2322);
\draw[gpiImprove] (4.0495,1.2322) -- (4.3047,0.9224);
\draw[gpiEval] (4.3047,0.9224) -- (4.3940,1.1584);
\draw[gpiImprove] (4.3940,1.1584) -- (4.5458,0.9741);
\end{scope}
\begin{scope}[xshift=6.9cm]
\node[font=\footnotesize\bfseries] at (2.45,2.65) {Partial policy evaluation};
\draw[SessionMuted,line width=.6pt] (0,2.1) -- (4.9,1.05) -- (0,0);
\node[font=\scriptsize,anchor=south,fill=SessionPaper,inner sep=1pt] at (2.1,1.77) {\(V=V^\sigma\)};
\node[font=\scriptsize,anchor=north] at (2.25,.34) {\(\sigma\) greedy w.r.t. \(V\)};
\fill[SessionCharcoal] (4.9,1.05) circle (1.1pt);
\node[font=\scriptsize,anchor=west] at (4.96,1.05) {\(V^*,\sigma^*\)};
\fill[SessionCharcoal] (0.15,0.65) circle (1pt);
\draw[gpiEval] (0.1500,0.6500) -- (0.8625,1.3458);
\draw[gpiImprove] (0.8625,1.3458) -- (2.0738,0.4444);
\draw[gpiEval] (2.0738,0.4444) -- (2.4977,1.0606);
\draw[gpiImprove] (2.4977,1.0606) -- (3.2184,0.6897);
\draw[gpiEval] (3.2184,0.6897) -- (3.4706,1.0563);
\draw[gpiImprove] (3.4706,1.0563) -- (3.8994,0.8356);
\draw[gpiEval] (3.8994,0.8356) -- (4.0495,1.0538);
\draw[gpiImprove] (4.0495,1.0538) -- (4.3047,0.9224);
\draw[gpiEval] (4.3047,0.9224) -- (4.3940,1.0522);
\draw[gpiImprove] (4.3940,1.0522) -- (4.5458,0.9741);
\end{scope}
\end{tikzpicture}
\caption{Full and partial policy evaluation within GPI (schematic).}
\end{figure}
\vspace{-.15cm}
{\footnotesize
\textcolor{CourseVariableColor}{\textbf{Policy evaluation}} moves toward \(\modelvar{V=V^\sigma}\).
\textcolor{SessionTeal}{\textbf{Policy improvement}} makes \(\modelvar{\sigma}\) greedy with respect to \(\modelvar{V}\).
\par\smallskip
At the intersection, both conditions hold: \(\modelvar{V=T_\sigma V=TV}\), so \(\modelvar{V=V^*}\) and \(\modelvar{\sigma}\) is optimal.
\par\smallskip
\textcolor{SessionMuted}{Concept adapted from Sutton and Barto (2018), Section 4.6, pp. 86--87.}}
\end{frame}

\begin{frame}{Modified Policy Iteration}
\small
In GPI, the evaluation arrows may stop short of full consistency. \term{Modified policy iteration} implements this using a fixed number \(\modelvar{m\geq1}\) of policy evaluation updates between policy improvements.
\begin{enumerate}
\item \textbf{Policy improvement.} Choose a policy greedy with respect to the current candidate \(\modelvar{V_k}\):
\[\sigma_{k+1}(s)\in\argmax_{a\in\Aset(s)}\{r(s,a)+\beta\E[V_k(s')\mid s,a]\}.\]
\item \textbf{Policy evaluation.} Hold \(\modelvar{\sigma_{k+1}}\) fixed. Set \(\modelvar{V_{k,0}=V_k}\), then compute
\[V_{k,j+1}=T_{\sigma_{k+1}}V_{k,j},\qquad j=0,\ldots,m-1.\]
\item Set \(\modelvar{V_{k+1}=V_{k,m}}\) and repeat.
\end{enumerate}
\par\smallskip
\(\modelvar{V_k}\) is an approximate value, so an unchanged greedy policy alone does not certify optimality. Check the Bellman residual \(\modelvar{\|TV_k-V_k\|_\infty}\), as in VFI.
\end{frame}

\begin{frame}{Choosing the Amount of Policy Evaluation}
\small
\textbf{One policy evaluation update: VFI.} Since \(\modelvar{\sigma_{k+1}}\) is greedy with respect to \(\modelvar{V_k}\),
\[T_{\sigma_{k+1}}V_k=TV_k.\]
Setting \(\modelvar{m=1}\) therefore gives \(\modelvar{V_{k+1}=TV_k}\).
\par\medskip
\textbf{Several policy evaluation updates: modified policy iteration.} Extra updates reuse the chosen actions. They avoid a new maximization while refining the value estimate.
\par\medskip
\textbf{Full policy evaluation: PFI.} Continuing evaluation to its fixed point gives \(\modelvar{V^{\sigma_{k+1}}}\); the next policy improvement uses that policy's value.
\par\medskip
For a finite discounted MDP with bounded rewards, these exact updates converge to \(\modelvar{V^*}\), including any fixed \(\modelvar{m\geq1}\).
\par\smallskip
The choice balances the cost of policy evaluation against the number of policy improvements.
\end{frame}

\begin{frame}{Policy Iteration with Continuous States}
\small
For the consumption--saving model, the state \(\modelvar{s}\) and saving choice \(\modelvar{a}\) are continuous:
\[
V^\sigma(s)=u(s-\sigma(s))+\beta\E[V^\sigma(\sigma(s)+y)],
\]
\[
\widetilde\sigma(s)\in\argmax_{0\leq a\leq\min\{s,M\}}
\{u(s-a)+\beta\E[V^\sigma(a+y)]\}.
\]
Here \(\modelvar{y}\) is income, and \(\modelvar{M}\) is the saving limit.
\begin{itemize}
\item \textbf{Represent the value function:} store values on a grid and interpolate between grid points.
\item \textbf{Compute expectations:} approximate the income integral using quadrature or simulation.
\item \textbf{Choose actions:} search an action grid or use continuous optimization.
\end{itemize}
\par\smallskip
Solving the resulting finite MDP accurately does not eliminate its grid and integration errors.
\end{frame}

% BEGIN MANAGED READINGS: week4_pfi
\begin{frame}{References and Further Reading}
\coursereading{1}
  {Richard S. Sutton and Andrew G. Barto (2018), Reinforcement Learning: An Introduction. 2nd ed. MIT Press. [Book]}
  {https://www.andrew.cmu.edu/course/10-703/textbook/BartoSutton.pdf}
  {\begin{itemize}
\item \href{https://www.andrew.cmu.edu/course/10-703/textbook/BartoSutton.pdf\#page=95}{Chapter 4: Dynamic Programming}
\end{itemize}}
\coursereading{2}
  {Thomas J. Sargent and John Stachurski, QuantEcon. [Online lectures]}
  {https://quantecon.org/lectures/}
  {\begin{itemize}
\item \href{https://tools-techniques.quantecon.org/discrete_dp.html\#policy-function-iteration}{Discrete State Dynamic Programming — Policy Function Iteration}
\item \href{https://tools-techniques.quantecon.org/discrete_dp.html\#modified-policy-function-iteration}{Discrete State Dynamic Programming — Modified Policy Function Iteration}
\end{itemize}}
\end{frame}
% END MANAGED READINGS: week4_pfi
\end{document}
